
% Options for packages loaded elsewhere
\PassOptionsToPackage{unicode}{hyperref}
\PassOptionsToPackage{hyphens}{url}
\PassOptionsToPackage{dvipsnames,svgnames,x11names}{xcolor}
%
\documentclass[
 11pt,
]{article}
\usepackage{amsmath,amssymb}
\usepackage{lmodern}
\usepackage{iftex}
\ifPDFTeX
 \usepackage[T1]{fontenc}
 \usepackage[utf8]{inputenc}
 \usepackage{textcomp} % provide euro and other symbols
\else % if luatex or xetex
 \usepackage{fontspec}
 \usepackage{textcomp} % text symbols (incl. quotes)
 \defaultfontfeatures{Scale=MatchLowercase}
 \defaultfontfeatures[\rmfamily]{Ligatures=TeX,Scale=1}
\fi
% Use upquote if available, for straight quotes in verbatim environments
\IfFileExists{upquote.sty}{\usepackage{upquote}}{}
\IfFileExists{microtype.sty}{% use microtype if available
 \usepackage[]{microtype}
 \UseMicrotypeSet[protrusion]{basicmath} % disable protrusion for tt fonts
}{}
\makeatletter
\@ifundefined{KOMAClassName}{% if non-KOMA class
 \IfFileExists{parskip.sty}{%
 \usepackage{parskip}
 }{% else
 \setlength{\parindent}{0pt}
 \setlength{\parskip}{6pt plus 2pt minus 1pt}}
}{% if KOMA class
 \KOMAoptions{parskip=half}}
\makeatother
\usepackage{xcolor}
\usepackage{newunicodechar}
% Map common Unicode math symbols appearing in text to LaTeX (avoid missing glyphs).
\newunicodechar{Ω}{\ensuremath{\Omega}}
\newunicodechar{Π}{\ensuremath{\Pi}}
\newunicodechar{Δ}{\ensuremath{\Delta}}
\newunicodechar{Θ}{\ensuremath{\Theta}}
\newunicodechar{μ}{\ensuremath{\mu}}
\newunicodechar{λ}{\ensuremath{\lambda}}
\newunicodechar{κ}{\ensuremath{\kappa}}
\newunicodechar{ρ}{\ensuremath{\rho}}
\newunicodechar{η}{\ensuremath{\eta}}
\newunicodechar{ε}{\ensuremath{\varepsilon}}
\newunicodechar{δ}{\ensuremath{\delta}}
\newunicodechar{α}{\ensuremath{\alpha}}
\newunicodechar{≤}{\ensuremath{\le}}
\newunicodechar{≥}{\ensuremath{\ge}}
\newunicodechar{∈}{\ensuremath{\in}}
\newunicodechar{≈}{\ensuremath{\approx}}

\newunicodechar{−}{\ensuremath{-}}
\newunicodechar{⇒}{\ensuremath{\Rightarrow}}
\newunicodechar{⇔}{\ensuremath{\Leftrightarrow}}
\newunicodechar{⋆}{\ensuremath{\star}}
\newunicodechar{∧}{\ensuremath{\wedge}}
\newunicodechar{□}{\ensuremath{\square}}
\newunicodechar{∃}{\ensuremath{\exists}}
\newunicodechar{∏}{\ensuremath{\prod}}
\newunicodechar{∩}{\ensuremath{\cap}}
\newunicodechar{∪}{\ensuremath{\cup}}
\newunicodechar{∅}{\ensuremath{\emptyset}}
\newunicodechar{⊆}{\ensuremath{\subseteq}}
\newunicodechar{∉}{\ensuremath{\notin}}
\newunicodechar{≠}{\ensuremath{\neq}}
\newunicodechar{∼}{\ensuremath{\sim}}
\newunicodechar{≲}{\ensuremath{\lesssim}}
\newunicodechar{≳}{\ensuremath{\gtrsim}}
\newunicodechar{≫}{\ensuremath{\gg}}
\newunicodechar{σ}{\ensuremath{\sigma}}
\newunicodechar{β}{\ensuremath{\beta}}
\newunicodechar{τ}{\ensuremath{\tau}}
\newunicodechar{γ}{\ensuremath{\gamma}}
\newunicodechar{ℓ}{\ensuremath{\ell}}
\newunicodechar{≪}{\ensuremath{\ll}}
\newunicodechar{⋃}{\ensuremath{\bigcup}}
\newunicodechar{′}{'}

\IfFileExists{xurl.sty}{\usepackage{xurl}}{} % add URL line breaks if available
\IfFileExists{bookmark.sty}{\usepackage{bookmark}}{\usepackage{hyperref}}
\hypersetup{
 colorlinks=true,
 linkcolor={blue},
 filecolor={Maroon},
 citecolor={Blue},
 urlcolor={blue},
 pdfcreator={LaTeX via pandoc}}
\urlstyle{same} % disable monospaced font for URLs
\usepackage[margin=1in]{geometry}
\setlength{\emergencystretch}{3em} % prevent overfull lines
\providecommand{\tightlist}{%
 \setlength{\itemsep}{0pt}\setlength{\parskip}{0pt}}
\setcounter{secnumdepth}{2} % number sections for referee readability
\setcounter{tocdepth}{2}

\title{Many-Testing Lower Bounds for Low-Latency Privacy on Hierarchical Overlays}
\author{}
\date{2026-01-20\\Last updated: 2026-02-24}

\begin{document}

\maketitle

\begin{abstract}
We formalize a global-surveillance failure mode for low-latency routed systems:
\textbf{many-testing}. A global passive adversary (GPA) can evaluate a large
family of candidate witnesses (prefix regions, cuts, bins, edges) and win by
selecting whichever witness is accidentally ``clean'' under cover traffic.
Even when any fixed witness is rare, a best-of-$B$ scan forces
\emph{logarithmic} overhead to suppress the witness across the entire family.

We give sharp emptiness / low-count thresholds (Poissonized and fixed-total),
and a dependence-robust theorem that lower-bounds scan success using a packing
number under overlap. We instantiate the framework on hierarchical
(prefix-partitioned) overlays to obtain a canonical cover-rate floor under a
deadline window of length $L$: for typical $B\approx n$ candidate prefix
regions, the cover rate must satisfy $\lambda \gtrsim \frac{1}{L}\ln(n/\varepsilon)$
to make scan distinguishers fail with probability at most $\varepsilon$.
Finally, we show that when edges are observable, hiding the \emph{neighbor}
choice requires per-edge (not just per-node) equalization, yielding an
$\Omega(\log^2 n/L)$-type floor in classic DHT scaling.
\end{abstract}

\tableofcontents
\newpage

\section{Introduction and
contributions}\label{introduction-and-contributions}

The guiding phenomenon is simple:

\begin{itemize}
\tightlist
\item
 A GPA can test many candidate ``places'' where a protocol's behavior
 might become visible (prefix regions, edges, bins, committees).
\item
 If a crisp rare event (like an empty baseline count) serves as a
 witness, then the GPA's best-of-many advantage forces logarithmic
 overhead to suppress that witness everywhere.
\end{itemize}

This paper is intentionally proof-oriented. It is not a new protocol; it
is a \textbf{lower-bound lens} that explains why several natural
defenses ``hit a floor'' in low-latency regimes.

Main contributions (informal):
\begin{itemize}\tightlist
\item \textbf{Sharp many-testing thresholds.} Exact emptiness/low-count scan bounds under
Poissonization and fixed totals, capturing the GPA's best-of-$B$ advantage.
\item \textbf{Overlap-robust scans.} A packing-number theorem that lower-bounds many-testing
success even when candidate witnesses overlap.
\item \textbf{Hierarchical-overlay instantiation.} Applying the framework to prefix regions
yields the canonical low-latency cover floor $\lambda \gtrsim \frac{1}{L}\ln(n/\varepsilon)$
(in typical $B\approx n$ regimes).
\item \textbf{Edge observability floor.} If the transcript reveals which neighbor edge is used,
then per-edge equalization is necessary, giving an $\Omega(\log^2 n/L)$-type floor in
classic DHT scaling.
\end{itemize}

Attack-as-distinguisher viewpoint. The many-testing phenomenon is not
just an existential bound: it yields an explicit transcript-level
distinguisher. the low-count scan scans a tested family (e.g., all depth-k
prefix regions) and flags a lookup iff some region has an unusually
small observed total. Lemma 4.2 gives closed-form false-positive and
true-positive bounds, and Proposition 4.1 translates them into the
canonical cover-rate floor by substituting μ = n λ L / B.

\subsubsection{Threat model and
observables}\label{threat-model-and-observables}

We consider a \textbf{global passive adversary (GPA)} who observes a
transcript over a fixed time window of length \textbf{L} rounds (or an
equivalent bounded-latency window).

\textbf{Test family.} Fix a finite family of ``places'' that can be
scanned in the transcript: I = \{1,2,\ldots,B\}. Each i∈I corresponds to
some region / cut / bin / committee / edge-event predicate S\_i (e.g., a
prefix region in a hierarchical overlay). The GPA observes a nonnegative
integer statistic C\_i for each i, typically ``how many messages crossed
/ were directed to / were attributed to S\_i'' inside the window.

\textbf{Null model (baseline cover).} Under pure cover traffic, we model
C\_i as approximately Poisson with mean μ\_i, often specialized to a
common mean μ. The many-testing theorems in §2 use exactly Poisson (or
fixed-total occupancy) to get sharp thresholds; §3 explains how overlap
or local dependence affects constants rather than the log-law.

\textbf{Alternative model (forced entry).} Under a correctness
requirement, the protocol forces bounded ``signature touches'' in a
small number of tested regions early in the window. Formally, we assume
there is a target-dependent index i⋆ such that, with failure probability
at most η, an additional nonnegative increment
F\_\{i⋆\}∈\{1,2,\ldots,m\} is added to the baseline count: C\_\{i⋆\} =
X\_\{i⋆\} + F\_\{i⋆\}, where X\_\{i⋆\} is the baseline component. The
adversary does \textbf{not} need to separate X and F; all distinguishers
we analyze use only the observable totals \{C\_i\}.

This formulation captures the common situation where ``baseline = 0'' is
unobservable but ``total is unusually small'' (e.g., C\_\{i⋆\}≤m) is
observable and still yields a witness.

\section{Setting and observables}\label{setting-and-observables}

We consider a GPA observing metadata at the network layer. The weakest
model we analyze is count-only per observed unit (e.g., per directed
edge per round). Stronger models (packet traces with timing/order) can
only help the adversary.

\subsubsection{Tested families and transcript
counts}\label{tested-families-and-transcript-counts}

Let I index a family of tests. Each test i∈I corresponds to a region /
cut / bin / edge whose \textbf{total observed count} Ci over a fixed
observation window of L rounds is available to the GPA (e.g., ``\#
packets entering prefix region S'', ``\# uses of edge e'', ``\# contacts
to bin (c,b)'').

For analysis it is useful to decompose Ci = Bi + Wi, where Bi is
\textbf{baseline cover} (workload-independent) traffic and Wi is
\textbf{workload-dependent} traffic (real lookups, responses, etc.). The
GPA observes Ci but does not (by design) label which packets are
baseline vs workload.

A canonical witness in many-testing arguments is \textbf{baseline
emptiness}: Ei := \{Bi = 0\}. Ei is generally not directly observable
from Ci. However, correctness can force that some target-dependent test
i⋆ receives a small, bounded number of workload packets within the
window.

Define the \textbf{observable low-count witness} (parameter m≥1): Li(m)
:= \{1 ≤ Ci ≤ m\}. This is measurable from the transcript whenever Ci
is.

\textbf{Lemma 1.1 (Observable low-count witnesses dominate baseline
emptiness).} Fix i and write Ci = Bi + Wi with Bi,Wi nonnegative
integers. 1) If Wi ≥ 1 then Li(m) implies Bi ≤ m−1. 2) If 1 ≤ Wi ≤ m
then Ei ∧ \{Wi ≥ 1\} implies Li(m). 3) If Wi = w0 deterministically for
some 1≤w0≤m, then \{Ci = w0\} is equivalent to Ei and is therefore an
observable surrogate for baseline emptiness.

Proof: (1) If Ci ≤ m and Wi ≥ 1 then Bi = Ci−Wi ≤ m−1. (2) If Bi=0 and
1≤Wi≤m then Ci=Wi∈{[}1,m{]}. (3) If Wi=w0 then Ci=w0 iff Bi=0. □

We will use Theorem 2.3 below to quantify how large μ must be to
suppress low-count witnesses Li(m) across many tests.

\section{Emptiness thresholds}\label{emptiness-thresholds}

\paragraph{Setup.}
Consider \(B\) candidate witnesses/tests. Under baseline-only traffic, each
test produces an independent count \(C_i \sim \mathrm{Pois}(\mu)\).
For an integer \(m\ge 0\), define the low-tail probability
\[
p_m(\mu) := \Pr[\mathrm{Pois}(\mu)\le m]
= e^{-\mu}\sum_{j=0}^m \frac{\mu^j}{j!}.
\]

\textbf{Theorem 2.3 (Many-testing threshold for low counts).}
Let \(E_m := \{\min_{i\in[B]} C_i \le m\}\) be the \emph{low-count scan}
event. Then
\[
\Pr[E_m] \;=\; 1-(1-p_m(\mu))^B.
\]
Consequently, requiring \(\Pr[E_m]\le \varepsilon\) forces (and is achieved
up to lower-order terms by)
\[
\mu \;\ge\; \ln\!\Big(\tfrac{B}{\varepsilon}\Big)
\;+\; m\ln\ln\!\Big(\tfrac{B}{\varepsilon}\Big)
\;-\; \ln(m!) \;-\; O(1).
\]
In particular, the emptiness case \(m=0\) yields the sharp logarithmic floor
\(\mu \ge \ln(B/\varepsilon)-O(1)\).

\emph{Proof sketch.}
The exact formula is by independence.
For the threshold, for \(\mu\ge 2m\) the Poisson low tail satisfies the
two-sided bound
\[
\frac{1}{2}\,e^{-\mu}\frac{\mu^m}{m!}
\;\le\; p_m(\mu) \;\le\;
2\,e^{-\mu}\frac{\mu^m}{m!},
\]
so in the rare-witness regime \(\Pr[E_m]\approx B\,p_m(\mu)\) and inverting
\(B\,p_m(\mu)\lesssim \varepsilon\) gives the stated expression.

\textbf{Corollary 2.4 (\(r\)-hit low-count scans).}
Let \(Z := |\{i: C_i\le m\}|\sim \mathrm{Bin}(B,p_m(\mu))\) and define
\(E_{m,r}:=\{Z\ge r\}\).
In the rare-witness regime \(t:=B\,p_m(\mu)\le 1\) and \(B\ge 2r\),
\[
\Pr[E_{m,r}] \;=\; \Theta\!\Big(\frac{t^r}{r!}\Big),
\]
so \(\Pr[E_{m,r}]\le \varepsilon\) implies \(B\,p_m(\mu)\lesssim (\varepsilon r!)^{1/r}\),
equivalently
\[
\mu \;\ge\; \ln\!\Big(\tfrac{B}{\varepsilon^{1/r}}\Big)
\;+\; m\ln\ln\!\Big(\tfrac{B}{\varepsilon^{1/r}}\Big)
\;-\; \ln(m!) \;-\; O(1).
\]

\paragraph{Remark (fixed-total occupancy).}
If a window contains exactly \(M\) baseline events thrown uniformly into \(B\) tests,
the same logarithmic floor holds with \(\mu \approx M/B\); e.g., suppressing
emptiness with probability at most \(\varepsilon\) requires
\(M \ge B\ln(B/\varepsilon)-O(B)\).
\section{Dependency-robust many-testing via packing and local
dependence}\label{dependency-robust-many-testing-via-packing-and-local-dependence}

Many natural test families overlap (nested prefix cuts, overlapping
windows). We isolate an effective number of independent tests.

\subsubsection{Poisson-atom baseline and packing
number}\label{poisson-atom-baseline-and-packing-number}

Let Ω be a set of independent Poisson atoms: for each a∈Ω, Na
\textasciitilde{} Pois(μa) independent. For each test i choose an atom
set Ai⊆Ω and define Bi := sum\_\{a∈Ai\} Na, λi :=
E{[}Bi{]}=sum\_\{a∈Ai\} μa. Then Bi\textasciitilde Pois(λi). If Ai∩Aj=∅
then Bi and Bj are independent.

Define the disjoint packing number: Π(\{Ai\}) :=
max\{\textbar J\textbar{} : J⊆I and Ai∩Aj=∅ for all i≠j in J\}.

\textbf{Theorem 3.1 (Packing-number many-testing bound).} Let λmax =
max\_i λi. Then P{[}exists i: Bi=0{]} ≥ 1 − (1 −
e\textsuperscript{\{−λmax\})}\{Π\}. In particular, enforcing P{[}exists
i: Bi=0{]} ≤ ε implies λmax ≥ ln(Π/ε) + O(1).

Proof: choose a disjoint subfamily J of size Π; on J, P{[}Bi=0{]} ≥
e\^{}\{−λmax\} and events are independent. □

\paragraph{Remark (tools and examples).} The packing reduction can be phrased via dependency graphs and Chen--Stein Poisson approximation; worked overlap examples and a broader ``minimal test'' discussion are omitted here for brevity.

\section{Prefix partitions and the Ω(ln n / L) cover floor (toy
baseline)}\label{prefix-partitions-and-the-ux3c9ln-n-l-cover-floor-toy-baseline}

This section instantiates the many-testing lower bound on a stylized hierarchical
overlay in a \emph{three-step pipeline}: (i) identify a large family of
target-dependent regions that can be tested from the transcript; (ii) argue that
a correct lookup forces a small, \emph{bounded} amount of workload traffic into
\emph{one} such region; (iii) apply the many-testing threshold to conclude a
cover-rate floor.

\paragraph{Step 1: a size-$B$ prefix family.}
Nodes have $b$-bit identifiers. For a target $t$, define the depth-$k$ prefix
region $S_k(t)$ as the set of nodes whose first $k$ bits match $t$. The family
$\mathcal{F}_k := \{S_k(u)\}_{u\in\{0,1\}^k}$ has size $B=2^k$, and for random IDs
each region has size about $n/B$.

\paragraph{Toy baseline (UCR).}
Uniform Constant Rate (UCR) baseline: each node emits $\lambda$ baseline packets
per round to a uniform random destination. For a region $S$,
\[
\mu(S) := \mathbb{E}[\text{baseline ingress into }S\text{ over }L\text{ rounds}]
= \lambda L\,|S|\,(n-|S|)/(n-1),
\]
so for balanced depth-$k$ regions ($|S|\approx n/B$ and $B\ll n$),
$\mu(S)\approx \lambda L n/B$. Applying Theorem~2.1 across $B$ disjoint tests
already suggests the canonical scale: to suppress best-of-$B$ emptiness with
failure probability $\le \varepsilon$, one needs $\mu \gtrsim \ln(B/\varepsilon)$,
i.e.\ $\lambda \gtrsim (B/(nL))\ln(B/\varepsilon)$, and in the deepest regime
$B\approx n$ this becomes $\lambda \gtrsim (1/L)\ln(n/\varepsilon)$.

\paragraph{Step 2: a transcript-measurable forced-entry witness.}
Let $C(S)$ be the \textbf{total observed} number of packets entering $S$ in the
$L$-round window, and write
$C(S)=B(S)+W(S)$ where $B(S)$ is baseline cover traffic and $W(S)$ is
workload-dependent traffic.

\textbf{Assumption 4.1 (Forced entry with bounded signature).}
For the true target $t$, with probability at least $1-\eta$ there exists a
depth-$k$ region $S_k(t)$ such that $1\le W(S_k(t))\le m$ for some small fixed
$m$, and the GPA observes only the totals $\{C(S)\}$ (so events like
$\{C(S_k(t))\le m\}$ are transcript-measurable).

\paragraph{Step 3: one scan distinguisher.}
\textbf{Definition 4.1 (Low-count scan).}
Given the transcript, compute the tested family $\mathcal{F}$ (e.g., all
depth-$k$ prefix regions) and output ``forced entry'' iff
$\min_{S\in\mathcal{F}} C(S)\le m$.

\textbf{Lemma 4.2 (Low-count scan advantage).}
Let $\mathcal{F}$ be a family of $B$ tests whose baseline ingress counts are
effectively independent with $B(S)\sim \mathrm{Pois}(\mu)$ under $H_0$
(baseline only). Under $H_1$ (forced entry) assume Assumption~4.1 holds.
Let $p_j(\mu):=\Pr[\mathrm{Pois}(\mu)\le j]$. Then the low-count scan has
\[
P_0 \;=\; 1-(1-p_m(\mu))^B,\qquad
P_1 \;\ge\; (1-\eta)\,p_{m-1}(\mu).
\]
\emph{Proof sketch.} Under $H_0$, the scan triggers iff at least one of the $B$
tests has $C(S)\le m$. Under $H_1$, on the forced-entry event there is a region
with $1\le W(S)\le m$, and whenever $B(S)\le m-1$ we have $C(S)=B(S)+W(S)\le m$.

\textbf{Proposition 4.1 (Canonical cover floor from forced entry via low counts).}
Assume there are $B$ effectively independent depth-$k$ regions whose baseline
ingress counts are i.i.d.\ $\mathrm{Pois}(\mu)$ with $\mu \approx n\lambda L/B$ in
UCR. Suppose Assumption~4.1 holds and we require that a GPA's scan succeeds with
probability at most $\varepsilon$ under baseline-only traffic. Then necessarily
\[
\mu \;\ge\; \ln(B/\varepsilon)+O(1),
\]
and more sharply (for fixed $m$) $\mu \ge \ln(B/\varepsilon)+ m\ln\ln(B/\varepsilon)
-\ln(m!)-O(1)$ as in Theorem~2.3. Substituting $\mu\approx n\lambda L/B$ yields
the cover-rate floor $\lambda \gtrsim (B/(nL))\ln(B/\varepsilon)$, and in the
deepest regime $B\approx n$ this reduces to $\lambda \gtrsim (1/L)\ln(n/\varepsilon)$.

\paragraph{Why Assumption 4.1 is natural (idealized).}
In a stylized greedy step, if a node chooses the best of $m$ candidate contacts
whose IDs are uniform among those matching the current prefix, the additional
prefix gain $G$ satisfies $\Pr[G\ge r]=1-(1-2^{-r})^m$ and $\mathbb{E}[G]\in[\log_2 m,\log_2 m+1]$.
This ``best-of-$m$ buys prefix bits'' heuristic explains why hierarchical
lookups tend to inject a small, target-dependent signature into progressively
deeper prefix regions.
\section{Edge observability and per-edge scheduling
floors}\label{edge-observability-and-per-edge-scheduling-floors}

The many-testing arguments above focus on \emph{counts over a family of
tests}. A different limitation applies when the GPA observes traffic
\emph{per edge} (or per neighbor) and the defense attempts to hide which
edges are ``active'' by per-edge padding.

\subsubsection{A basic per-edge latency--rate
tradeoff}\label{a-basic-per-edge-latencyrate-tradeoff}

Consider a route that traverses h directed edges. Suppose each edge e
emits cover packets according to a schedule with rate ρ (packets per
round). In the strongest ``edge-observable'' model, the GPA sees
per-edge emission times (or equivalently, counts in each round per
edge).

A minimal requirement for low latency is: a real payload packet should
traverse the h-hop path within L rounds with high probability. If the
schedule is the only way to move packets (i.e., payload rides on
scheduled emissions), then each hop incurs an expected waiting time at
least 1/ρ.

\textbf{Lemma 7.1 (Edge scheduling floor; rate grows with path length).}
In the per-edge scheduled-emission model above, any scheme that
guarantees expected end-to-end latency at most L on an h-hop path must
satisfy ρ ≥ h/L for every edge on that path. Consequently, for a node of
out-degree Δ, the per-node scheduled-slot rate must be at least Δρ = Ω(Δ
h / L).

\emph{Proof.} Under any stationary per-edge schedule of rate ρ, the
expected time until the next emission on an edge is at least 1/ρ (e.g.,
by renewal theory; for a Poisson clock it is exactly 1/ρ, and for
deterministic periodic schedules it is also 1/ρ). A packet traversing h
edges therefore has expected latency at least h·(1/ρ). Requiring h/ρ ≤ L
yields ρ ≥ h/L. Summing over Δ outgoing edges gives the per-node rate
lower bound. □

\subsubsection{Classic DHT scaling: Ω(log² n /
L)}\label{classic-dht-scaling-ux3c9loguxb2-n-l}

In many overlays, h=Θ(log n) hops and Δ=Θ(log n) outgoing neighbors are
typical scales (e.g., logarithmic-degree routing tables and logarithmic
route length).

\textbf{Corollary 7.2 (Edge-observability floor in logarithmic
overlays).} If Δ=Θ(log n) and h=Θ(log n), then any per-edge schedule
that maintains expected latency ≤ L requires per-node scheduled-slot
rate Ω(log² n / L).

This edge-observability floor is orthogonal to the many-testing and 
 floors: it applies even when there is only a single tested
region, as long as the GPA can observe traffic per edge.

\section{Synthesis: what the lower bounds force}\label{synthesis}

The results above can be read as a short design checklist for low-latency
metadata-hiding lookups.

\begin{itemize}\tightlist
\item \textbf{Best-of-many scans impose logarithmic mass.} If a GPA can scan an effective family
of size $B$, then suppressing emptiness/low-count witnesses across the family forces baseline
mass $\mu=\Omega(\ln(B/\varepsilon))$ in the relevant regions.
\item \textbf{Hierarchical structure turns $B$ into $n$.} In prefix-partitioned overlays, natural
candidate families have size comparable to the population ($B\approx n$), yielding the canonical
cover floor $\lambda \gtrsim \frac{1}{L}\ln(n/\varepsilon)$ in deadline $L$.
\item \textbf{Observable edges require per-edge scheduling.} When the transcript reveals the neighbor
edge used, hiding the choice requires per-edge (not just per-node) equalization, producing additional
degree-dependent floors (e.g., $\Omega(\log^2 n/L)$ under classic DHT scaling).
\end{itemize}

These statements are intentionally protocol-agnostic: they describe what \emph{any} defense must
pay in order to make the corresponding scan distinguishers fail.

\section{Related work and
positioning}\label{related-work-and-positioning}

Our goal is not to propose a full protocol but to isolate
\emph{proof-friendly, transcript-measurable witnesses} that a global
passive adversary (GPA) can test at scale, and to show that suppressing
those witnesses imposes logarithmic (or worse) baseline overhead in
low-latency regimes.

\textbf{Overlay motivation.} The routing picture is inspired by
XOR-metric DHT overlays such as Kademlia {[}MM02{]} and by anonymous DHT
lookup constructions that attempt to hide query targets using dummy
traffic and/or split paths (e.g., Octopus {[}WB12{]}). These systems
naturally induce large families of ``regions'' (prefix buckets, routing
cuts) that can be scanned for rare events such as \emph{no contact} or
\emph{few contacts}.

\textbf{Low-latency anonymity motivation.} On the communication-systems
side, low-latency anonymity work highlights the tension between rate
equalization and global traffic analysis, e.g., Loopix-style Poisson
mixing with cover traffic {[}PH+17{]}, round-based private communication
systems such as Pung {[}A+16{]}, and distributed metadata-private
messaging such as Stadium {[}T+17{]}. Our contribution is complementary:
even if a design equalizes \emph{marginals} (per-node or per-bin), a GPA
can often test \emph{many} regions/bins/edges and win by selecting
whichever produces a crisp rare event.

\textbf{Probabilistic backbone.} The many-testing phenomena we use are
classical in Poisson/occupancy asymptotics and Poisson approximation: Le
Cam's Poisson-binomial approximation {[}LC60{]} (and expositions such as
{[}Ste94{]}), dependency-graph Chen--Stein bounds in the style of
Arratia--Goldstein--Gordon {[}AGG89{]}, and large-deviation Poisson
approximation heuristics (e.g., {[}Jan90{]}). We package these tools in
a way that is easy to instantiate for hierarchical overlays and
low-latency privacy designs.

\subsubsection{Positioning}\label{positioning}

A recurring theme is that ``equalize the marginal'' (per-node rate,
per-bin usage, per-edge slots) is not enough when the GPA can test
\emph{many} places and look for a crisp, rare event. Our contribution is
to turn that intuition into referee-checkable inequalities with minimal
protocol detail (isolated in explicit assumptions like Assumption 4.1).

\section{Limitations and open problems}\label{limitations}

\begin{itemize}\tightlist
\item \textbf{Transcript model.} We focus on GPA-visible count statistics over a fixed window $L$.
Extending the scan framework to richer observables (timing, correlation, queueing) is important.
\item \textbf{Effective test size.} The parameter $B$ is application-dependent; formalizing \emph{which}
candidate families a real adversary can scan (and how dependent they are) remains case-specific.
\item \textbf{Protocol specificity.} Our instantiations are stylized (prefix regions, edge choices).
Sharper constants and protocol-specific bounds require matching the witness family to an implementation.
\item \textbf{Achievability.} Designing defenses that match these floors under realistic constraints
(bandwidth caps, churn, partial observability) is still largely open.
\end{itemize}

\subsection{References}\label{references}

{[}AGG89{]} R. Arratia, L. Goldstein, and L. Gordon. ``Two moments
suffice for Poisson approximations: the Chen--Stein method.''
\emph{Annals of Probability} 17(1): 9--25, 1989.

{[}A+16{]} S. Angel, T. Chen, A. Laine, and others. ``Unobservable
Communication over Fully Untrusted Infrastructure.'' \emph{OSDI} 2016.
(System name: Pung.)

{[}B+20{]} L. Barman, I. Dacosta, M. Zamani, E. Zhai, A. Pyrgelis, B.
Ford, J. Feigenbaum, and J.-P. Hubaux. ``PriFi: Low-Latency Anonymity
for Organizational Networks.'' \emph{Proceedings on Privacy Enhancing
Technologies} 2020(4): 20--43. DOI: 10.2478/popets-2020-0059.

{[}GRS12{]} A. Ghosh, T. Roughgarden, and M. Sundararajan. ``Universally
Utility-Maximizing Privacy Mechanisms.'' \emph{SIAM Journal on
Computing} 41(6): 1673--1693, 2012. DOI: 10.1137/09076828X.

{[}GV16{]} Q. Geng and P. Viswanath. ``The Optimal Noise-Adding
Mechanism in Differential Privacy.'' \emph{IEEE Transactions on
Information Theory} 62(2): 925--951, 2016.

{[}Jan90{]} S. Janson. ``Poisson approximation for large deviations.''
\emph{Random Structures \& Algorithms} 1: 221--229, 1990.

{[}LC60{]} L. Le Cam. ``An approximation theorem for the Poisson
binomial distribution.'' \emph{Pacific Journal of Mathematics} 10(4):
1181--1197, 1960.

{[}LGZ19{]} D. Lazar, Y. Gilad, and N. Zeldovich. ``Yodel: Strong
Metadata Security for Voice Calls.'' \emph{SOSP} 2019. DOI:
10.1145/3341301.3359644.

{[}MM02{]} P. Maymounkov and D. Mazières. ``Kademlia: A Peer-to-Peer
Information System Based on the XOR Metric.'' In \emph{Peer-to-Peer
Systems (IPTPS 2002)}, LNCS 2429, pp.~53--65. Springer, 2002. DOI:
10.1007/3-540-45748-8\_5.

{[}PH+17{]} A. M. Piotrowska, J. Hayes, T. Elahi, S. Meiser, and G.
Danezis. ``The Loopix Anonymity System.'' \emph{USENIX Security
Symposium} 2017.

{[}Ste94{]} J. M. Steele. ``Le Cam's Inequality and Poisson
Approximations.'' \emph{The American Mathematical Monthly} 101(1):
48--54, 1994.

{[}T+17{]} N. Tyagi, Y. Gilad, D. Leung, M. Zaharia, and N. Zeldovich.
``Stadium: A Distributed Metadata-Private Messaging System.''
\emph{SOSP} 2017.

{[}WB12{]} Q. Wang and N. Borisov. ``Octopus: A Secure and Anonymous DHT
Lookup.'' \emph{IEEE ICDCS} 2012, pp.~325--334. DOI:
10.1109/ICDCS.2012.78.

\begin{itemize}
\tightlist
\item
 R. Motwani and P. Raghavan. \emph{Randomized Algorithms}. Cambridge
 University Press, 1995. (Balls-into-bins and coupon collector.)
\item
 M. Mitzenmacher and E. Upfal. \emph{Probability and Computing}.
 Cambridge University Press, 2005. (Concentration and occupancy
 bounds.)
\end{itemize}

\end{document}